\documentclass[11pt]{article}

\usepackage[T1]{fontenc}
\usepackage[utf8]{inputenc}
\usepackage{amsmath}
\usepackage{amssymb}
\usepackage{amsthm}
\usepackage{booktabs}
\usepackage{geometry}
\geometry{a4paper, margin=25mm}
\usepackage{hyperref}
\emergencystretch=3em

\newtheorem{theorem}{Theorem}
\newtheorem{lemma}{Lemma}
\newtheorem{proposition}{Proposition}
\newtheorem{definition}{Definition}
\newtheorem{hypothesis}{Hypothesis}
\renewcommand{\thehypothesis}{H\arabic{hypothesis}}
\newtheorem{remark}{Remark}

\newcommand{\E}{\mathbb{E}}
\newcommand{\Prob}{\mathbb{P}}
\newcommand{\Var}{\operatorname{Var}}
\newcommand{\cost}{\mathrm{cost}}
\newcommand{\dto}{\xrightarrow{d}}
\newcommand{\lbl}[1]{\texttt{#1}}

\title{A batched observable for growing variance\\
\large bringing $\sigma_i^2\to\infty$ back under the article's CLT}
\author{\texttt{loglog\_validation} project notes}
\date{2026-11-13}

\begin{document}
\maketitle

\begin{abstract}
When the normalised variance of an observable grows like
$\sigma_i^2\sim\sigma_*^2 i^{2\nu}$, two hypotheses of the article's
Theorem~\lbl{thm:CLT} fail: Assumption~\lbl{sigma\_converges}, and the
uniform $2+\delta$ moment bound of Assumption~\lbl{2\_plus\_delta\_moments},
which already bounds $\sigma_i^2$. This note replaces $Y_i$ by the mean
$Z_i$ of $K_i=\lceil c\,i^{2\nu}\rceil$ independent copies. $Z_i$ has the same
mean, hence the same $\gamma$, $a_1$ and $\omega_1$, and its normalised
variance converges. Under one condition on the growth of the
$(2+\delta)$-th moment it also satisfies the uniform moment bound, so every
assumption of Theorem~\lbl{thm:CLT} holds and the article's CLT, Wilson bound
and planner apply to $Z_i$ unchanged. The condition holds at equality for an
indicator of a rare event (Leath's $2-\tau$), follows from the standard
cluster-size scaling ansatz for the susceptibility, and is what the
measurements show for $\gamma_{\rm susc}$ in $d=3,4,5$. The cost per sample
becomes $\propto i^{d+2\nu}$, so the allocation is \lbl{def:opt} with
$d\to d+2\nu$, the same rate as the direct heteroscedastic analysis of
\texttt{heteroscedastic\_allocation.tex}. At a fixed budget, batching costs
$12$--$14\%$ in standard deviation in the $d=5$ example. That is the price of
staying inside the theorem.
Nothing here is in the article; it is an extension for the user to review.
\end{abstract}

\tableofcontents

% =================================================================
\section{What the CLT uses, and what breaks}

Labels refer to \texttt{../article\_writting/main.tex}. For an observable
$Y_i$ with $\E Y_i>0$ write $\xi_i=Y_i/\E Y_i$ and $\sigma_i^2=\Var\xi_i$.
Theorem~\lbl{thm:CLT} assumes:

\begin{enumerate}
  \item[(A1)] the $J$-order expansion
  $\E Y_i=a_0i^\gamma\exp\bigl(a_1i^{-\omega_1}+\dots+\phi_J(i)i^{-\omega_J}\bigr)$
  (Assumption~\lbl{firstAssumption-}\allowbreak\lbl{J-order-expansion}), with $\phi_J$
  bounded (Assumption~\lbl{phi\_is\_unif\_bounded});
  \item[(A2)] $\E Y_i>0$ (Assumption~\lbl{Y\_positive});
  \item[(A3)] $\E|\xi_i|^{2+\delta}\le M$ for all $i$, some $\delta>0$
  (Assumption~\lbl{2\_plus\_delta\_moments});
  \item[(A4)] $\sigma_i^2\to\sigma_\infty^2\in(0,\infty)$
  (Assumption~\lbl{sigma\_converges}).
\end{enumerate}

\noindent In the proof (Appendix, \lbl{subsec:moments}--\lbl{subsec:lyapunov}):
\begin{itemize}
  \item (A3) gives the Lyapunov bound of Lemma~\lbl{lem:lyapunov},
  $\sum\E|X_{k,\ell}|^{2+\delta}/s_{n,m}^{2+\delta}\le
  C\,M'/(\sigma_{\min}^{2+\delta}(nm)^{\delta/2})$, with
  $M'=(M^{1/(2+\delta)}+1)^{2+\delta}$ (Lemma~\lbl{lem:moments}\,(c)).
  \item (A4) gives $\sigma_{\min}$ bounded away from $0$ and the
  deterministic normalisation of Lemma~\lbl{lem:S-clt}.
  \item $\sigma_{\max}^2<\infty$, which (A3) implies, enters
  Lemma~\lbl{lem:Q-prob}: $\Prob(\mathcal G^c)\le4m\sigma_{\max}^2/n$ and the
  remainder $Q_{n,m}$ is $O(\sigma_{\max}^2/n)$.
\end{itemize}

\noindent Since $\Var\xi_i\le\E\xi_i^2\le(\E|\xi_i|^{2+\delta})^{2/(2+\delta)}$,
(A3) bounds $\sigma_i^2$ by $M^{2/(2+\delta)}$. A growing $\sigma_i$ therefore
violates (A3) as well as (A4). The last item is where it hurts in practice:
$4m\sigma_{\max}^2/n$ small is the condition that $\sigma/\sqrt n$ be small
at the top scale, the guard of \texttt{heteroscedastic\_allocation.tex}, which
the $d=5$ run of experiment~60 broke.

% =================================================================
\section{Hypotheses on the raw observable}

\begin{hypothesis}[Regular variance growth]\label{H:var}
There are $\sigma_*>0$ and $\nu\ge0$ such that
$\sigma_i^2=\sigma_*^2\,i^{2\nu}\,(1+\varepsilon_i)$ with $\varepsilon_i\to0$.
\end{hypothesis}

\begin{hypothesis}[Moment growth]\label{H:mom}
For the same $\delta>0$ as in (A3), the normalised absolute moment
\[
  \mu_i:=\frac{\E|\xi_i-1|^{2+\delta}}{\sigma_i^{2+\delta}}
\]
satisfies $\mu_i\le C_\mu\,i^{\nu\delta}$ for all $i$.
\end{hypothesis}

\noindent \ref{H:var} with $\nu=0$ is (A4). \ref{H:mom} allows the shape of the
tail to get heavier with $i$, at a rate tied to $\nu$; for $\nu=0$ it says
$\E|\xi_i-1|^{2+\delta}$ is bounded, i.e.\ (A3). Section~\ref{sec:H2} explains
why $\mu_i\asymp i^{\nu\delta}$, the boundary case, is the natural one for
rare-event observables. (A1) and (A2) are kept as they are.

% =================================================================
\section{The batched observable}

\begin{definition}\label{def:batch}
Fix $c>0$ and let $K_i=\lceil c\,i^{2\nu}\rceil$. One sample of the batched
observable at scale $i$ is
\[
  Z_i=\frac1{K_i}\sum_{r=1}^{K_i}Y_i^{(r)},
\]
with $Y_i^{(1)},\dots,Y_i^{(K_i)}$ independent copies of $Y_i$, fresh for
each sample. Write $\zeta_i=Z_i/\E Z_i=\frac1{K_i}\sum_r\xi_i^{(r)}$.
\end{definition}

\noindent A long run then draws $n$ samples of $Z_i$ at each scale of the
window, exactly as the article prescribes ($n_k\equiv n$); in raw samples
that is $nK_i$ at scale $i$.

\begin{proposition}[Mean: nothing changes]\label{prop:mean}
$\E Z_i=\E Y_i$. Hence $Z_i$ satisfies (A1) and (A2) with the same
$a_0,\gamma,a_1,\omega_1,\dots,\phi_J$.
\end{proposition}

\begin{proof}
Linearity of expectation.
\end{proof}

\noindent In particular, the correction to scaling is still $\omega_1$, not
$\omega_1+\nu$: no rescaling of $Y_i$ by a deterministic factor can change its
coefficient of variation, so the only way to tame $\sigma_i$ is to spend more
randomness per observation, and that shows up in the cost
(Section~\ref{sec:alloc}), not in the expansion of the mean.

\begin{proposition}[Variance converges]\label{prop:var}
Under \ref{H:var}, $\Var\zeta_i=\sigma_i^2/K_i\to\sigma_*^2/c=:\sigma_{Z,\infty}^2\in(0,\infty)$.
So $Z_i$ satisfies (A4).
\end{proposition}

\begin{proof}
The copies are independent, so $\Var\zeta_i=\sigma_i^2/K_i$. Since
$c\,i^{2\nu}\le K_i<c\,i^{2\nu}+1$, $K_i/(c\,i^{2\nu})\to1$ when $\nu>0$
(and $K_i=\lceil c\rceil$ is constant when $\nu=0$, where the claim is (A4)).
Then $\sigma_i^2/K_i=\sigma_*^2(1+\varepsilon_i)\,c\,i^{2\nu}/(c\,K_i)\to\sigma_*^2/c$.
\end{proof}

% -----------------------------------------------------------------
\subsection{The \texorpdfstring{$2+\delta$}{2+delta} moments}

The tool is Rosenthal's inequality: for $p\ge2$ there is $C_p<\infty$ such that,
for i.i.d.\ centred $X_1,\dots,X_K$ with $\E|X_1|^p<\infty$,
\begin{equation}
  \E\Bigl|\sum_{r=1}^KX_r\Bigr|^p\le C_p\Bigl[\bigl(K\,\E X_1^2\bigr)^{p/2}+K\,\E|X_1|^p\Bigr].
  \label{eq:rosenthal}
\end{equation}
The first term is what a Gaussian sum would give; the second is the
contribution of a single large summand.

\begin{proposition}[Uniform $2+\delta$ moment]\label{prop:mom}
Under \ref{H:var} and \ref{H:mom}, with $p=2+\delta$,
\[
  \E|\zeta_i-1|^{2+\delta}\le C_p\Bigl[\Bigl(\frac{\sigma_i^2}{K_i}\Bigr)^{1+\delta/2}
  +\frac{\mu_i\,\sigma_i^{2+\delta}}{K_i^{1+\delta}}\Bigr]
  \le C_p\Bigl[\bigl(\sigma_{Z,\infty}^2\bigr)^{1+\delta/2}
  +C_\mu\,\frac{\sigma_*^{2+\delta}}{c^{1+\delta}}\Bigr](1+o(1)),
\]
which is bounded in $i$. Consequently $\E|\zeta_i|^{2+\delta}\le M_Z$ for all
$i$, with $M_Z^{1/(2+\delta)}=1+\sup_i\|\zeta_i-1\|_{2+\delta}$, and $Z_i$
satisfies (A3).
\end{proposition}

\begin{proof}
Apply \eqref{eq:rosenthal} to $X_r=\xi_i^{(r)}-1$ and divide by $K_i^{p}$:
\[
  \E|\zeta_i-1|^p\le C_p\Bigl[\frac{(K_i\sigma_i^2)^{p/2}}{K_i^p}
  +\frac{K_i\,\E|\xi_i-1|^p}{K_i^p}\Bigr]
  =C_p\Bigl[\Bigl(\frac{\sigma_i^2}{K_i}\Bigr)^{p/2}+\frac{\mu_i\sigma_i^{p}}{K_i^{p-1}}\Bigr].
\]
The first term converges to $(\sigma_*^2/c)^{1+\delta/2}$ by
Proposition~\ref{prop:var}. In the second, with $K_i\ge c\,i^{2\nu}$ and
\ref{H:var}, \ref{H:mom},
\[
  \frac{\mu_i\sigma_i^{2+\delta}}{K_i^{1+\delta}}
  \le C_\mu i^{\nu\delta}\,
  \frac{\sigma_*^{2+\delta}(1+\varepsilon_i)^{1+\delta/2}\,i^{\nu(2+\delta)}}{c^{1+\delta}\,i^{2\nu(1+\delta)}}
  =C_\mu\frac{\sigma_*^{2+\delta}}{c^{1+\delta}}(1+\varepsilon_i)^{1+\delta/2}\,
   i^{\nu\delta+2\nu+\nu\delta-2\nu-2\nu\delta},
\]
and the exponent is $0$. Minkowski's inequality in $\mathcal L^{2+\delta}$,
$\|\zeta_i\|_{2+\delta}\le1+\|\zeta_i-1\|_{2+\delta}$, gives the last claim,
exactly as in Lemma~\lbl{lem:moments}\,(c) read backwards.
\end{proof}

\begin{remark}[\ref{H:mom} is the sharp condition for this argument]
The exponent bookkeeping shows the second Rosenthal term behaves like
$\mu_i\,i^{-\nu\delta}$. If $\mu_i$ grew faster than $i^{\nu\delta}$, the
term would diverge, and it cannot be dropped: for a sum dominated by one
large summand, \eqref{eq:rosenthal} is attained up to $C_p$. Batching tames
the \emph{variance} by $K_i$ but a single extreme copy only by $K_i^{-(1+\delta)}$
relative to its own $(2+\delta)$-th moment, and \ref{H:mom} is exactly what
makes these two match.
\end{remark}

\begin{remark}[$\delta\le1$ from $\delta=1$]
$\mathcal L^p$ norms are non-decreasing in $p$ on a probability space, so a
uniform bound on $\E|\zeta_i|^{3}$ gives one on $\E|\zeta_i|^{2+\delta}$ for
every $\delta\in(0,1]$. The repo measures $\delta=1$
(\texttt{tools/summary.py}, \texttt{abs\_moment} $=\E|\xi|^{3}$), which is
therefore enough for the CLT.
\end{remark}

\begin{remark}[What Lyapunov needs is weaker]
The Lyapunov ratio of Lemma~\lbl{lem:lyapunov} only has to tend to $0$, not to
be bounded by a constant over $(nm)^{\delta/2}$. Redoing its first display
with the per-scale moments instead of their supremum,
\[
  \frac{1}{s_{n,m}^{2+\delta}}\sum_{k,\ell}\E|X_{k,\ell}|^{2+\delta}
  \le\frac{6^\delta}{12^{\delta/2}\,\sigma_{\min}^{2+\delta}}
  \cdot\frac{\max_{k}\E|\zeta_{\rho^k}-1|^{2+\delta}}{(nm)^{\delta/2}},
\]
so a moment that grows slowly with the window is allowed as long as
$n$ outgrows it. (A3) is a sufficient condition, which is why
Proposition~\ref{prop:mom} aims for it.
\end{remark}

% -----------------------------------------------------------------
\subsection{The CLT for the batched observable}

\begin{theorem}[Theorem~\lbl{thm:CLT} for $Z_i$]\label{thm:batch}
Let $Y_i$ satisfy (A1), (A2), \ref{H:var} and \ref{H:mom}, and let $\hat\gamma$
be the log-log estimator \lbl{eq:hatbeta} computed from $n$ samples of the
batched observable $Z_{\rho^k}$ of Definition~\ref{def:batch} at each scale
$k=m_0+1,\dots,m_0+m$. If $m/n\to0$ and $m_0\to\infty$,
\[
  \sqrt{n\,m(m^2-1)}\Bigl(\hat\gamma-\gamma-\frac{R(m,m_0)}{\log\rho}\Bigr)
  \dto\mathcal N\Bigl(0,\frac{12\,\sigma_*^2}{c\,\log^2\rho}\Bigr),
\]
with $R(m,m_0)$ the finite-size bias \lbl{eq:bias} of $Y_i$ itself. The Wilson
bound (Theorem~\lbl{cor:wilson}) holds with $\sigma_\infty^2=\sigma_*^2/c$ and
$\sigma_{\max}^2=\sup_{k\ge m_0+1}\sigma_{\rho^k}^2/K_{\rho^k}$.
\end{theorem}

\begin{proof}
By Propositions~\ref{prop:mean}, \ref{prop:var} and~\ref{prop:mom}, $Z_i$
satisfies (A1)--(A4) with the same expansion and with
$\sigma_\infty^2=\sigma_*^2/c$. Theorem~\lbl{thm:CLT} applied to $Z_i$ is the
claim.
\end{proof}

\paragraph{Two consequences inside the proof.}
\begin{itemize}
  \item \emph{The remainder.} Lemma~\lbl{lem:Q-prob} now reads
  $\Prob(\mathcal G^c)\le4m\,\sigma_{Z,\max}^2/n$ with $\sigma_{Z,\max}$ bounded:
  the article's $m/n\to0$ suffices, and the guard on $\sigma/\sqrt n$ at the top
  of the window is no longer an extra condition. For the raw observable the same
  bound reads $4m\sigma_*^2\rho^{2\nu(m_0+m)}/n$, which needs $n$ to outgrow the
  top scale's variance.
  \item \emph{The log bias.} To second order
  $\E\hat\beta-\beta-R(m,m_0)\approx-\frac1{2n}\sum_kw_{k,m}\sigma_{\rho^k}^2/K_{\rho^k}$.
  With $\sigma^2/K$ equal at every scale up to $o(1)$ and $\sum_kw_{k,m}=0$, the
  sum is $o(1/n)$: the term $b_{\rm LOG}$ of the heteroscedastic note cancels
  as it does in the article.
\end{itemize}

% =================================================================
\section{Why \texorpdfstring{\ref{H:mom}}{H2} holds at equality for rare-event observables}\label{sec:H2}

\paragraph{An indicator: Leath's $2-\tau$.}
The model \texttt{percolation\_tau\_leath} has
$Y_s=\mathbf 1\{|C(0)|\ge s\}$, mean $Q(s)\sim q_0s^{2-\tau}$. Then
$\sigma_s^2=(1-Q)/Q\sim s^{\tau-2}$, so $\nu=(\tau-2)/2$, and
\[
  \E|\xi_s-1|^{2+\delta}=Q\Bigl(\frac{1-Q}{Q}\Bigr)^{2+\delta}+(1-Q)
  \sim Q^{-(1+\delta)},
  \qquad
  \mu_s\sim\frac{Q^{-(1+\delta)}}{Q^{-(1+\delta/2)}}=Q^{-\delta/2}\sim s^{\nu\delta}.
\]
\ref{H:mom} holds with equality. The batched observable is
$\zeta_s=\mathrm{Bin}(K_s,Q)/(K_sQ)$; with $K_sQ\to\lambda:=c\,q_0$ it converges
in distribution, with all moments, to $\mathrm{Poisson}(\lambda)/\lambda$, and
$\sigma_{Z,\infty}^2=1/\lambda$. This is the model's existing
\texttt{batch\_factor}$\cdot s^{\texttt{batch\_exponent}}$ with the exponent
$2\nu$. Two practical notes: $\Prob(\zeta_s=0)\to e^{-\lambda}$, so $\lambda$
of a few units keeps $\overline Z>0$ with overwhelming probability for any
reasonable $n$; and $\nu=(\tau-2)/2$ is a function of the answer, so it must be
estimated (Section~\ref{sec:nu}), never set from a known $\tau$ (ground rule 4;
the model's docstring already says so).

\paragraph{The susceptibility, under the scaling ansatz.}
For the size of the origin's cluster below $p_c$, the standard ansatz
$n_s\simeq s^{-\tau}f(s/s_\xi)$, $s_\xi\sim\varepsilon^{-1/\sigma}$, with
$\Prob(|C|=s\mid0\text{ open})\propto s\,n_s$, gives
$\E|C|^k\asymp s_\xi^{\,k+2-\tau}$ for every $k\ge1$ (the sum
$\sum_s s^{k+1}n_s$ is dominated by $s\approx s_\xi$ since $k+2>\tau$; $k=1$ is
$S\asymp s_\xi^{3-\tau}$, i.e.\ $\gamma=(3-\tau)/\sigma$). For
$\xi=|C|/\E|C|$,
\begin{gather*}
  \sigma^2\asymp\frac{\E|C|^2}{(\E|C|)^2}\asymp s_\xi^{\,(4-\tau)-2(3-\tau)}=s_\xi^{\,\tau-2},
  \\
  \E|\xi|^{2+\delta}\asymp\frac{\E|C|^{2+\delta}}{(\E|C|)^{2+\delta}}
  \asymp s_\xi^{\,(4+\delta-\tau)-(2+\delta)(3-\tau)}=s_\xi^{\,(1+\delta)(\tau-2)}.
\end{gather*}
Hence $\E|\xi|^{2+\delta}\asymp\sigma^{2+2\delta}$ and
$\mu\asymp\sigma^{\delta}\asymp i^{\nu\delta}$: \ref{H:mom} at equality, for
every $\delta$, as for the indicator. The shared structure is that both
observables are carried by rare samples of size comparable to the cutoff,
which is what makes $\E|\xi|^{2+\delta}$ grow like $\sigma^{2+2\delta}$
rather than like $\sigma^{2+\delta}$. The repo's explore model sums over seeds on
an axis rather than taking one cluster, which changes $\nu$ (it is $0.21$--$0.39$,
not the single-cluster value), so this is a guide, not a proof for that
observable.

\paragraph{The measurements.}
From the pilots and long runs (\texttt{abs\_moment} and \texttt{cv} per scale),
fitting the growth exponent of $\mu_i=\E|\xi_i|^3/\sigma_i^3$ against $\nu$
($\delta=1$):

\begin{center}
\begin{tabular}{@{}lrrrr@{}}
\toprule
study & scales & $\mu_i$, first $\to$ last & growth of $\mu_i$ & $\nu\delta$\\
\midrule
$\gamma_{\rm susc}$, $d=3$ (exp.~51) & $4\ldots2048$ & $8.0\to25.8$ & $0.191$ & $0.210$\\
$\gamma_{\rm susc}$, $d=4$ (exp.~54) & $2\ldots512$ & $9.9\to56.1$ & $0.316$ & $0.308$\\
$\gamma_{\rm susc}$, $d=5$ (exp.~60) & $4\ldots131072$ & $15.3\to660$ & $0.395$ & $0.394$\\
$2-\tau$, $d=3$ (exp.~56) & $16\ldots262144$ & $6.3\to3.6$ & $-0.041$ & $0.138$\\
$2-\tau$, $d=5$ (exp.~57) & $16\ldots16777216$ & $3.4\to25.5$ & $0.163$ & $0.215$\\
\bottomrule
\end{tabular}
\end{center}

\noindent The susceptibility sits on the boundary $\mu_i\asymp i^{\nu}$ in all
three dimensions, as the ansatz says. The Leath rows are below it because
$Q$ is not yet small on the lower scales (in $d=3$ the local $\nu$ is itself
still falling). Two cautions: these are sample moments of heavy-tailed
variables, which are biased low and more so at the top scales; and only
$\delta=1$ is measured. \ref{H:mom} at equality is the marginal case of
Proposition~\ref{prop:mom}: the bound holds, with no room to spare.

% =================================================================
\section{Cost and allocation}\label{sec:alloc}

\paragraph{Cost.} A sample of $Z_i$ costs $K_i$ samples of $Y_i$:
\[
  \cost_Z(i)=K_i\,\cost(i)\approx c\,i^{2\nu}\,(\kappa_0+\kappa_1i^{d}).
\]
With the article's power cost (Assumption~\lbl{cost\_is\_power\_law}),
$\cost_Z(i)\propto i^{d+2\nu}$, and Lemma~\lbl{lem:budget} and
Definition~\lbl{def:opt} apply to $Z_i$ with $d':=d+2\nu$:
\[
  n=C_\rho'B^{\theta_1},\quad m_0=\theta_2\log_\rho B,\qquad
  \theta_1=\frac{2\omega_1}{d+2\nu+2\omega_1},\quad
  \theta_2=\frac{1}{d+2\nu+2\omega_1},
\]
and the error decays as $B^{-\omega_1/(d+2\nu+2\omega_1)}$. In raw samples
of $Y$ the scale $\rho^k$ gets $nK_{\rho^k}\propto B^{\theta_1+2\nu\theta_2}$ at
the bottom of the window, i.e.\ $B^{(2\omega_1+2\nu)/(d+2\nu+2\omega_1)}$: the
same exponents as the direct heteroscedastic analysis. The two routes agree on
the rate; they differ in how samples are spread inside the window (below).

\paragraph{The planner needs no change.} Step~2 of the article
(\lbl{eq:n-of-m0}, \lbl{eq:bSE}, \lbl{eq:bFS}, \lbl{eq:m0-star}) is used as is,
with $\widehat\cost_Z$ in place of $\widehat\cost$ and
$\hat\sigma^2_{Z,\infty}=\hat\sigma_*^2/c$ in place of $\hat\sigma^2_\infty$.
The constant $c$ drops out: $n$ scales as $1/c$ and $\sigma_{Z,\infty}^2$ as
$1/c$, so \lbl{eq:bSE} is unchanged. $c$ only sets granularity; since
$K_i=\lceil c\,i^{2\nu}\rceil$ is a rounded quantity, choose it so that $K$ is
at least $\sim10$ at the bottom of the window (relative rounding error
$\le1/K$), e.g.\ $K_i=\lceil K_0\,(i/i_0)^{2\nu}\rceil$ with $i_0$ the window's
smallest scale.

\paragraph{What staying inside the theorem costs.} At a fixed budget $B$ and
fixed window, with per-scale raw allocation $n_k$, $B\Var\hat\beta$ is, for
the three choices of $n_k$:
\begin{align*}
  \text{equal raw $n$ (the article on $Y$)}&:\ \Bigl(\sum_kw_k^2\sigma_k^2\Bigr)\Bigl(\sum_k\cost_k\Bigr),\\
  \text{batched, $n_k\propto\sigma_k^2$}&:\ \Bigl(\sum_kw_k^2\Bigr)\Bigl(\sum_k\sigma_k^2\cost_k\Bigr),\\
  \text{Neyman, $n_k\propto|w_k|\sigma_k/\sqrt{\cost_k}$}&:\ \Bigl(\sum_k|w_k|\sigma_k\sqrt{\cost_k}\Bigr)^2,
\end{align*}
and Cauchy--Schwarz makes the Neyman split the smallest of the three. On the
$d=5$ study of experiment~60 ($m=6$, measured $\sigma_k$ and cost per
sample), relative to equal raw $n$:

\begin{center}
\begin{tabular}{@{}lccc@{}}
\toprule
window & batched / equal (sd) & Neyman / equal (sd) & $K_{\rm top}/K_{\rm bottom}$\\
\midrule
$512\ldots16384$ & $1.123$ & $0.875$ & 17\\
$4096\ldots131072$ & $1.136$ & $0.870$ & 10\\
\bottomrule
\end{tabular}
\end{center}

\noindent Batching moves samples to the expensive top of the window, where
$w_k^2$ is no larger than at the bottom, so it loses $12$--$14\%$ in standard
deviation against equal $n$. What it buys: the article's theorems apply
(Theorem~\ref{thm:batch}); the OLS weights \lbl{eq:hatbeta} are optimal again,
since the scales are homoscedastic; $b_{\rm LOG}$ cancels; and $\sigma/\sqrt n$
is the same at every scale instead of largest at the top, which is where the
$d=5$ run's Gaussian approximation failed. Equal raw $n$ is cheaper only while
its top scale stays inside that approximation.

% =================================================================
\section{\texorpdfstring{$\nu$}{nu} must be estimated}\label{sec:nu}

$K_i$ needs $\nu$, and for Leath $\nu=(\tau-2)/2$ is a function of the
exponent being measured. Ground rule~4 rules out setting it from a known value.
The pilot supplies $\hat\nu$ by least squares of $\log\hat\sigma_i$ on $\log i$
(for Leath, of $\hat\sigma_i^2+1=1/\hat Q$, which is a pure power). Two facts
make this legitimate:

\begin{itemize}
  \item The pilot's samples are independent of the long run's (ground rule~2),
  so, conditionally on the pilot, $K_i=\lceil c\,i^{2\hat\nu}\rceil$ is a
  deterministic sequence and Theorem~\ref{thm:batch}'s proof goes through
  verbatim for the long run.
  \item With $\hat\nu\ne\nu$, $\Var\zeta_i\propto i^{2(\nu-\hat\nu)}$, i.e.\ the
  batched observable has a small residual growth $\nu-\hat\nu$. Strictly (A4)
  fails again, but over the window the variance changes by
  $\rho^{2m|\nu-\hat\nu|}$: with the $d=5$ pilot's $\pm0.006$ and $m=6$ that
  is $5\%$, which the planner can carry by the heteroscedastic formula of the
  companion note or simply ignore.
\end{itemize}

% =================================================================
\section{Open points}

\begin{enumerate}
  \item \ref{H:mom} is measured only at $\delta=1$, with sample moments that
  are biased low for heavy tails. For the susceptibility it follows from the
  scaling ansatz, which is itself a hypothesis; for the explore model's axis
  sum it is only measured.
  \item \ref{H:mom} holds at equality, so Proposition~\ref{prop:mom} gives a
  bound with no slack: a logarithmic correction to the moment growth would break
  uniform boundedness, though the weaker Lyapunov condition (Remark~3) would
  still hold for $n$ large enough.
  \item Raw samples per batched sample grow like $i^{2\nu}$; models that
  already average internally (the explore model's axis seeds, the torus
  sum of $d=2$) are a partial batching with a fixed exponent. Making that
  exponent equal to $2\hat\nu$ is the cleanest implementation.
\end{enumerate}

\end{document}
